\section{Harmonic defect, conditional savings, and the native wall}
\label{sec:tpc61-harmonic-defect}

The row ladder shows that one omitted small cofactor is more dangerous than
many omitted cofactors near the critical ceiling.  We now turn this fact
into a literal weighted estimate and audit its endpoint strength.

\subsection{A weighted harmonic certificate}

For each row and orbit time define the coefficient-free high weight
\begin{equation}
 \Omega_m(j)
 :=\sum_{\substack{w\in\cW_H(y)\\m(w)=m,\ j(w)=j}}
       |\beta_w|^2.
 \label{eq:tpc61-Omega-mj}
\end{equation}
Every summand at a fixed \((m,j)\) has the same recovered cofactor
\begin{equation}
 r_m(j):=\frac{mj+h_0}{\Pplus(mj+h_0)}
 \label{eq:tpc61-realized-row-cofactor}
\end{equation}
whenever \(\Pplus(mj+h_0)>y\).  Different opposite-row atoms remain
distinct in
\eqref{eq:tpc61-Omega-mj}.  Thus
\begin{equation}
 W_H(m)=\sum_j\Omega_m(j),
 \qquad
 W_M(m)=\sum_{\substack{j\\r_m(j)\notin\cI_m}}\Omega_m(j).
 \label{eq:tpc61-row-weight-orbit-form}
\end{equation}

For an active row put
\begin{equation}
 U_m:=\max_j\Omega_m(j),
 \qquad
 H_m^{\rm eff}:=\frac{W_H(m)}{U_m}.
 \label{eq:tpc61-effective-high-exposure}
\end{equation}
The second quantity is a weighted effective number of exposed orbit times;
it lies in \((0,\#\mathcal J_X]\).  Let
\begin{equation}
 E_m:=\{r\in\N:r\le N_+/y,\ r\notin\cI_m\}
 \label{eq:tpc61-missing-cofactor-set}
\end{equation}
and define its rowwise sieve-harmonic capacity by
\begin{equation}
 \mathfrak h_m(E_m)
 :=\frac m{\varphi(m)}
   \sum_{r\in E_m}
   \left(
    1+\frac{J/r}{\log(2+J/r)}
   \right).
 \label{eq:tpc61-harmonic-defect}
\end{equation}
Including unrealized \(r\)'s in \(E_m\) only enlarges this upper
certificate.

\begin{theorem}[Literal harmonic-defect bound]
\label{thm:tpc61-harmonic-defect-bound}
For every active row,
\begin{equation}
 \boxed{
 W_M(m)\ll U_m\mathfrak h_m(E_m),}
 \label{eq:tpc61-row-missing-harmonic-bound}
\end{equation}
and hence
\begin{equation}
 \boxed{
 \varepsilon_m
 \ll\frac{\mathfrak h_m(E_m)}{H_m^{\rm eff}}.}
 \label{eq:tpc61-row-defect-harmonic-bound}
\end{equation}
Consequently, if a number \(\epsilon_X\ge0\) satisfies
\begin{equation}
 \sup_{m:W_H(m)>0}
 \frac{\mathfrak h_m(E_m)}{H_m^{\rm eff}}
 \le\epsilon_X,
 \label{eq:tpc61-uniform-harmonic-hypothesis}
\end{equation}
then, with an absolute fixed-carrier constant,
\begin{equation}
 \cD_M(\zeta)\ll\epsilon_X\cD_H(\zeta)
 \label{eq:tpc61-uniform-harmonic-transfer}
\end{equation}
for every coefficient vector in the declared uniform class.
\end{theorem}

\begin{proof}
Group the second sum in \eqref{eq:tpc61-row-weight-orbit-form} by the
recovered cofactor.  For fixed \(r\), there are \(N_m(r)\) eligible orbit
times, and each aggregate weight is at most \(U_m\).  Therefore
\begin{equation}
 W_M(m)
 \le U_m\sum_{r\in E_m}N_m(r).
 \label{eq:tpc61-WM-before-BT}
\end{equation}
Apply \cref{prop:tpc61-brun-titchmarsh-row-bound} and use
\eqref{eq:tpc61-harmonic-defect}; this proves
\eqref{eq:tpc61-row-missing-harmonic-bound}.  Division by
\(W_H(m)=U_mH_m^{\rm eff}\) gives
\eqref{eq:tpc61-row-defect-harmonic-bound}.  The uniform conclusion follows
from \cref{thm:tpc61-rowwise-exposure-minimax}.
\end{proof}

The theorem is unconditional as an inequality, but a useful numerical
value of its right side requires two arithmetic inputs: completion of the
retained relation in harmonic capacity and a lower bound for the actual
effective high exposure.

\subsection{Dyadic exponent law}

At the critical endpoint,
\begin{equation}
 \theta_J:=\frac{133}{400},
 \qquad
 J=X^{\theta_J+o(1)},
 \qquad
 \frac{N_+}{y}=X^{\theta_J+o(1)}
 \label{eq:tpc61-critical-J-exponent}
\end{equation}
when \(y=C_{\rm vh}Q\) with the required fixed support constant
\citep{WangTPC59}.  Suppose a dyadic block
\(R=X^{\theta+o(1)}\), \(0\le\theta\le\theta_J\), satisfies uniformly
on the claimed rows
\begin{equation}
 |E_m\cap[R,2R]|\le X^{e(\theta)+o(1)},
 \qquad
 H_m^{\rm eff}\ge X^{\eta_H+o(1)}.
 \label{eq:tpc61-dyadic-defect-exposure-hypotheses}
\end{equation}
For this block write
\begin{align}
 \cW_{M,R}
 &:=\{w\in\cW_M:R\le r(w)<2R\},
 \label{eq:tpc61-missing-block-carrier}\\
 W_{M,R}(m)
 &:=\sum_{\substack{R\le r<2R\\r\notin\cI_m}}\nu_m(r),
 \qquad
 \cD_{M,R}(\zeta):=\sum_mW_{M,R}(m)|\zeta_m|^2.
 \label{eq:tpc61-missing-block-diagonal}
\end{align}

\begin{corollary}[Dyadic cofactor-saving exponent]
\label{cor:tpc61-dyadic-saving-exponent}
The contribution of this block to the row defect obeys
\begin{equation}
 \frac{W_{M,R}(m)}{W_H(m)}
 \le
 X^{-\sigma_{\rm cof}(\theta)+o(1)},
 \label{eq:tpc61-block-defect-saving}
\end{equation}
where
\begin{equation}
 \boxed{
 \sigma_{\rm cof}(\theta)
 :=\eta_H-e(\theta)-(\theta_J-\theta).}
 \label{eq:tpc61-general-cofactor-saving}
\end{equation}
In the balanced-exposure case \(\eta_H=\theta_J\), this simplifies to
\begin{equation}
 \boxed{
 \sigma_{\rm cof}(\theta)=\theta-e(\theta).}
 \label{eq:tpc61-balanced-cofactor-saving}
\end{equation}
The conclusion is useful only when the displayed saving is positive.
An \(O(\log X)\) dyadic union changes no fixed-power exponent.
\end{corollary}

\begin{proof}
By \cref{cor:tpc61-dyadic-missing-census}, the harmonic numerator from
the block has exponent at most
\(e(\theta)+\theta_J-\theta\).  Divide by the effective-exposure lower
bound in \eqref{eq:tpc61-dyadic-defect-exposure-hypotheses}.  Summing a
logarithmic number of nonnegative block bounds contributes \(X^{o(1)}\).
\end{proof}

\begin{example}[One small missing cofactor defeats global density]
\label{ex:tpc61-one-small-cofactor}
Suppose all integers up to the critical ceiling are retained except
\(r=1\).  Its unweighted global missing fraction is
\(J^{-1+o(1)}=X^{-133/400+o(1)}\), but the row ladder is
\begin{equation}
 p=mj+h_0.
 \label{eq:tpc61-r-one-ladder}
\end{equation}
The support capacity has \(J\) candidates and Brun--Titchmarsh removes
only a logarithm.  Hence this cardinal statement certifies no positive
fixed-power saving.  If the actual high weight on a row is supported on
the realized points of this ladder, then \(\varepsilon_m=1\).  The current
retained-relation interface contains no exposure hypothesis excluding that
possibility.
\end{example}

\begin{example}[Weight concentration defeats cofactor counting]
\label{ex:tpc61-weight-concentration}
Even if the missing set has much smaller cardinality than the represented
set, the exact measure can satisfy
\(\nu_m(\cI_m)=0<\nu_m(\cI_m^c)\) when all represented profile squares vanish
and one missing profile square is nonzero.  Less extremely, represented
weights may be arbitrarily smaller than a missing weight.  Thus cardinality
can replace \(\nu_m\) only after a two-sided physical-weight or effective-
exposure estimate has been proved on the same carrier.
\end{example}

These examples explain why the correct support statistic is harmonic and
why even harmonic smallness must be divided by actual exposure.  In
particular, a proof based only on this certificate must cover every
cofactor \(r<X^{\sigma-o(1)}\) exactly, or supply an independent
prime/weight saving on the omitted small-cofactor ladders, before claiming
an \(X^{-\sigma}\) missing-mass bound.

\subsection{The residual wedge and the direct-native closure}

The preceding exponent can now be compared with the grouping costs of
TPC--60.  Write
\begin{equation}
 \lambda_L:=\frac{99979}{210000}.
 \label{eq:tpc61-lambda-L}
\end{equation}
On a block \(R=X^{\theta+o(1)}\), the ambient residual-key degree from the
fixed-determinant ladder is
\begin{equation}
 d_{\rm res}(R)\le X^{\lambda_L-\theta+o(1)}.
 \label{eq:tpc61-residual-degree-exponent}
\end{equation}
In the balanced-exposure scenario with missing-cardinality exponent
\(e(\theta)\), the phase-blind residual mass--degree exponent is therefore
\begin{equation}
 \lambda_L-2\theta+e(\theta).
 \label{eq:tpc61-residual-product-exponent}
\end{equation}

At the critical block \(\theta=133/400\),
\begin{equation}
 \lambda_L-\theta
 =\frac{15077}{105000},
 \label{eq:tpc61-critical-residual-degree}
\end{equation}
whereas
\begin{equation}
 \boxed{
 2\theta-\lambda_L
 =\frac{39671}{210000}
 =0.1889095\ldots.}
 \label{eq:tpc61-residual-wedge}
\end{equation}
Thus the residual-key exponents themselves leave a genuine positive
window.  They do not close the physical map, because the latter also
erases the residual label \(s\).

To state the remaining requirement in the notation of the restricted-Gram
section, factor the missing map on one block as
\begin{equation}
 \cK_X
 \xrightarrow{\ L_{M,R}\ }
 \ell^2(\cW_{M,R})
 \xrightarrow{\ C_{{\rm res},R}\ }
 \ell^2(i,s)
 \xrightarrow{\ \mathsf E\ }
 \ell^2(i),
 \label{eq:tpc61-two-stage-missing-map}
\end{equation}
where \(T_{M,R}:=C_{{\rm res},R}L_{M,R}\) is the residual map on the block and
\(\mathsf E\) is the canonical residual-label erasure.  Define the
blockwise restricted erasure constant
\begin{equation}
 \beta_{M,R}
 :=\left\|
 \mathsf E\big|_{\Ran T_{M,R}}
 \right\|^2
 =X^{\mu_{\rm erase}+o(1)}.
 \label{eq:tpc61-restricted-erasure-exponent}
\end{equation}
This is the block analogue of \(\beta_M\) in
\cref{thm:tpc61-restricted-erasure-norm}; it is not the ambient number of
residual labels at one native output.  Likewise
\eqref{eq:tpc61-residual-degree-exponent} gives
\begin{equation}
 \kappa_{{\rm res},R}
 :=\sup_{\cD_{M,R}(\zeta)>0}
   \frac{\|T_{M,R}\zeta\|^2}{\cD_{M,R}(\zeta)}
 \le d_{\rm res}(R)
 \le X^{\lambda_L-\theta+o(1)}.
 \label{eq:tpc61-block-residual-kappa}
\end{equation}
If the represented reference satisfies
\begin{equation}
 \|A_X\zeta\|^2
 \ge X^{-\lambda_{\rm rep}+o(1)}\cD_H(\zeta),
 \label{eq:tpc61-represented-reference-loss}
\end{equation}
then the critical-block contribution has zero fixed-power reassembly cost
under the strict condition
\begin{equation}
 \boxed{
 e+\mu_{\rm erase}+\lambda_{\rm rep}
 <\frac{39671}{210000}.}
 \label{eq:tpc61-critical-residual-route-condition}
\end{equation}
Indeed, the two-stage inequality
\begin{equation}
 \|M_{X,R}\zeta\|^2
 \le\beta_{M,R}\kappa_{{\rm res},R}
       \cD_{M,R}(\zeta)
 \label{eq:tpc61-block-two-stage-upper}
\end{equation}
shows that the product of missing mass, residual degree, restricted
erasure, and represented-reference loss is then \(o(1)\).  The reverse
triangle certificate of TPC--60 yields the blockwise conclusion
\begin{equation}
 \frac{\|(A_X+M_{X,R})\zeta\|^2}{\|A_X\zeta\|^2}=1+o(1).
 \label{eq:tpc61-zero-cost-reassembly-consequence}
\end{equation}
For the full missing map, the same conclusion holds only if every other
nonempty dyadic block is empty or satisfies its analogous strict estimate
uniformly with a fixed positive margin.  Since there are \(O(\log X)\)
blocks, those estimates give
\begin{equation}
 \sum_R\|M_{X,R}\zeta\|=o(\|A_X\zeta\|),
 \label{eq:tpc61-all-blocks-small}
\end{equation}
and then \(M_X=\sum_RM_{X,R}\) yields the full version of
\eqref{eq:tpc61-zero-cost-reassembly-consequence}.  A favorable critical
block alone does not control missing mass on smaller cofactors.
The condition is uniform or distinguished according to the scope of
\eqref{eq:tpc61-represented-reference-loss} and the cofactor-exposure
hypotheses; the two scopes are not interchangeable.

By contrast, the only inherited direct native-degree bound has exponent
\begin{equation}
 \mu_{\rm native}=\frac{267}{400}.
 \label{eq:tpc61-crude-native-exponent}
\end{equation}
Even the best nonempty pure-cardinality certificate at the critical scale,
one missing cofactor with balanced exposure, saves at most \(133/400\).
Hence the direct-native phase-blind architecture has the fixed exponent
deficit
\begin{equation}
 \boxed{
 \frac{267}{400}-\frac{133}{400}
 =\frac{67}{200}>0.}
 \label{eq:tpc61-direct-native-deficit}
\end{equation}
It cannot prove reassembly by paying the full native degree.  This is a
route-specific closure, not a lower bound for the actual restricted norm
and not an endpoint loss to be inserted into the ledger.

\paragraph{Claim boundary.}
The exact row measure, cofactor ladder, Brun--Titchmarsh estimate, harmonic
certificate, and exponent audit are L0/L1 results at one fixed \(h_0\).
No estimate in this section proves that the actual retained relation has a
small harmonic complement, that every active row has polynomial high
exposure, or that \(\mu_{\rm erase}\) satisfies
\eqref{eq:tpc61-critical-residual-route-condition}.  Consequently there is
no L2 fixed-shift lower square, parity improvement, prime-pair asymptotic,
or twin-prime implication.  The positive conclusion is narrower and
falsifiable: the residual wedge is open, the coefficient-blind direct
native door is closed, and restricted M\"obius erasure is the next
arithmetic gate.
